\documentclass[11pt]{article}
\usepackage[T1]{fontenc}
\usepackage{lmodern,microtype}
\usepackage[margin=1in]{geometry}
\usepackage{amsmath,amssymb,amsthm,mathtools,booktabs}
\usepackage[hidelinks]{hyperref}
\usepackage{fancyhdr}
\pagestyle{fancy}\fancyhf{}
\fancyhead[L]{\small Complementary-subset dilations and hafnian bounds}
\fancyhead[R]{\small Research draft}\fancyfoot[C]{\thepage}
\setlength{\headheight}{14pt}\setlength{\emergencystretch}{2em}
\hypersetup{pdftitle={A complementary-subset dilation for sharp hafnian gradient inequalities},pdfauthor={No Free Collapse project - research draft}}
\newtheorem{theorem}{Theorem}
\newtheorem{lemma}[theorem]{Lemma}
\newtheorem{corollary}[theorem]{Corollary}
\theoremstyle{remark}\newtheorem{remark}[theorem]{Remark}
\DeclareMathOperator{\haf}{haf}\DeclareMathOperator{\tr}{tr}
\newcommand{\R}{\mathbb R}\newcommand{\op}{\mathrm{op}}
\title{A complementary-subset dilation for sharp\\hafnian gradient inequalities}
\author{No Free Collapse project\thanks{Proof draft. Human author names, affiliations, and the final assistance disclosure are to be confirmed before submission. Independent proof and priority review remain pending.}}
\date{12 September 2026}
\begin{document}\maketitle
\begin{abstract}
For a real symmetric matrix $A$ of order $2m$, we study the weighted Kneser
operator $\mathcal T_m(A)$ on $(m-1)$-subsets, whose disjoint-subset entries
are indexed by the two remaining vertices. An explicit complementary-subset
dilation gives the sharp bound
$\|\mathcal T_m(A)\|_{\op}\le(m+1)(\lambda_{\max}(A)-\lambda_{\min}(A))/4$
for every $m\ge2$. In order six, this proves the hafnian-gradient inequality
$q_2(A)\le d(A)^2q_1(A)/4$ and the energy inequality
$|\haf(A)|\le d(A)q_1(A)/6$, where $d(A)$ is the spectral diameter,
$q_1$ is off-diagonal squared energy, and $q_2$ is the sum of squared
complementary four-variable hafnians. We classify all equality cases for
the positive-semidefinite versions: apart from diagonal degeneracies, they
are scaled signed balanced rank-one projections, their rank-five complements,
and signed pair-block rank-three projections. The proof supplies explicit
positive-semidefinite certificates and avoids rank-by-rank projection arguments.
An exact-arithmetic implementation checks the dilation on matrix bases and
verifies rational instances. The distinct Boolean-cube-normalized hafnian
extremal problem is not resolved by these spectral bounds.
\end{abstract}
\noindent\textbf{Keywords:} hafnian; positive semidefinite matrix; Kneser operator;
spectral diameter; dilation; equality cases.

\section{Introduction and the six-variable question}
The hafnian of a symmetric matrix of even order is the sum of the products
of entries along its perfect matchings. For a four-set $S=\{i,j,k,l\}$,
\[
\haf(A[S])=a_{ij}a_{kl}+a_{ik}a_{jl}+a_{il}a_{jk}.
\]
For a real symmetric matrix of order six define
\begin{equation}\label{eq:q}
q_1(A)=\sum_{i<j}a_{ij}^2,\qquad
q_2(A)=\sum_{i<j}\haf(A_{\widehat i,\widehat j})^2.
\end{equation}
The second expression is the squared Euclidean norm of the gradient of the
hafnian when the off-diagonal entries are the independent variables.
The spectral inequality
\begin{equation}\label{eq:psdgoal}
q_2(A)\le\frac{\lambda_{\max}(A)^2}{4}q_1(A),\qquad A\succeq0,
\end{equation}
arose as an auxiliary problem in the No Free Collapse project. Earlier
arguments in that project established it on the real projection locus, using
separate rank-one, rank-two, and rank-three arguments and complementation.
Here we prove the stronger spectral-diameter bound for every real symmetric
matrix and identify the entire PSD equality locus.

Our main tool is an explicit signed incidence dilation. It applies in every
even order to a weighted Kneser operator, not just to the gradient problem in
order six. Neither the dilation nor its semidefinite consequence requires
numerical optimization. The equality proof uses orthogonality of the two
complementary-subset ranges and an elementary classification of the resulting
edge support.

Kneser adjacency, square-free polynomial spaces, and positive-semidefinite
compressions are established tools. Roos~\cite{roos} develops inequalities
for averaged squared subhafnians. Lindell and Staples~\cite{lindell} study
norm inequalities in zeon algebras. Kummer~\cite{kummer}, Lemma~3.5, uses an
unweighted Kneser adjacency operator in a square-free derivative isomorphism;
his main theorem concerns spectrahedral representations. Our statement
concerns a specified matrix-weighted operator, its sharp spectral-diameter
constant, and its six-variable equality cases. Section~\ref{sec:scope}
separates this statement from the original cube-constrained extremal problem.

\section{Complementary-subset dilation}
Let $n=2m$, $m\ge2$, and write $\mathcal E=\binom{[n]}{m-1}$.
For real symmetric $A\in\R^{n\times n}$, define the symmetric matrix
$\mathcal T_m(A)$ on $\R^{\mathcal E}$ by
\begin{equation}\label{eq:operator}
\mathcal T_m(A)_{E,F}=
\begin{cases}
a_{ij},& E\cap F=\varnothing,\quad [n]\setminus(E\cup F)=\{i,j\},\\
0,& E\cap F\ne\varnothing.
\end{cases}
\end{equation}
In particular, $\mathcal T_m(A+tI)=\mathcal T_m(A)$.
For every $m$-subset $R$, define the linear map $c_R:\R^{\mathcal E}\to\R^n$ by
\[
(c_R(v))_i=\begin{cases}v_{R\setminus\{i\}},&i\in R,\\0,&i\notin R.\end{cases}
\]
Choose one representative from each of the $N=\binom{2m}{m}/2$
complementary pairs $\{R,R^c\}$, and define $L_\pm$ by stacking
$c_R\pm c_{R^c}$ over these representatives.

\begin{lemma}[Exact dilation]\label{lem:dilation}
The incidence maps satisfy $L_\pm^TL_\pm=(m+1)I$, and
\begin{equation}\label{eq:dilation}
\mathcal T_m(A)=\frac14\bigl[
L_+^T(I_N\otimes A)L_+-L_-^T(I_N\otimes A)L_-\bigr].
\end{equation}
\end{lemma}
\begin{proof}
Each row of $L_\pm$ has exactly one nonzero entry, equal to $1$ or $-1$.
For a fixed $(m-1)$-subset $E$, its coordinate occurs once for each
$m$-subset containing $E$, hence $m+1$ times. This proves the Gram identities.
For one complementary pair, expansion gives
\[
(c_R+c_{R^c})^TA(c_R+c_{R^c})-
(c_R-c_{R^c})^TA(c_R-c_{R^c})=4c_R^TAc_{R^c}.
\]
For fixed disjoint $E,F$ with complement $\{i,j\}$, the monomial
$a_{ij}v_Ev_F$ occurs in two complementary partitions: the remaining vertices
can be assigned to $E,F$ in either order. Thus the sum of the cross terms is
$v^T\mathcal T_m(A)v$. Equality of these quadratic forms proves
\eqref{eq:dilation}.
\end{proof}

\begin{theorem}[Sharp spectral-diameter bound]\label{thm:main}
For every real symmetric $A$ of order $2m$, $m\ge2$,
\begin{equation}\label{eq:main}
\boxed{\|\mathcal T_m(A)\|_{\op}\le\frac{m+1}{4}
\bigl(\lambda_{\max}(A)-\lambda_{\min}(A)\bigr).}
\end{equation}
The coefficient is optimal for every $m\ge2$.
\end{theorem}
\begin{proof}
First suppose $0\preceq A\preceq\ell I$ and put $c=(m+1)/4$.
Lemma~\ref{lem:dilation} gives
\begin{align}
c\ell I-\mathcal T_m(A)
&=\tfrac14\bigl[L_+^T(I_N\otimes(\ell I-A))L_+
 +L_-^T(I_N\otimes A)L_-\bigr],\label{eq:minus}\\
c\ell I+\mathcal T_m(A)
&=\tfrac14\bigl[L_-^T(I_N\otimes(\ell I-A))L_-
 +L_+^T(I_N\otimes A)L_+\bigr].\label{eq:plus}
\end{align}
Both right sides are PSD, proving the norm bound $c\ell$.
Subtracting $\lambda_{\min}(A)I$ proves \eqref{eq:main} for general $A$.
For sharpness take $A=J_n/n$. Its spectral diameter is one. Each row of
$\mathcal T_m(A)$ has $\binom{m+1}{2}$ entries equal to $1/n$, so the
all-ones vector has eigenvalue $(m+1)/4$.
\end{proof}

\section{Hafnian inequalities in order six}
When $m=3$, the coordinates of $\mathcal T_3(A)$ are the fifteen edges of
$K_6$. Let $a=(a_{ij})_{i<j}$ and
$h=(\haf(A_{\widehat i,\widehat j}))_{i<j}$. Direct matching enumeration gives
\begin{equation}\label{eq:gradient}
\mathcal T_3(A)a=2h,\qquad a^Th=3\haf(A).
\end{equation}
The first factor counts the two orders of each complementary pair of edges;
the second counts each full matching once for each of its three edges.

\begin{corollary}\label{cor:gradient}
For a real symmetric matrix $A$ of order six, put
$d(A)=\lambda_{\max}(A)-\lambda_{\min}(A)$. Then
\begin{equation}\label{eq:corollaries}
\boxed{q_2(A)\le\frac{d(A)^2}{4}q_1(A),\qquad
|\haf(A)|\le\frac{d(A)}6q_1(A).}
\end{equation}
In particular \eqref{eq:psdgoal} holds for every PSD $A$.
\end{corollary}
\begin{proof}
Use $4q_2=\|\mathcal T_3(A)a\|^2\le d(A)^2\|a\|^2$ and
$6|\haf(A)|=|a^T\mathcal T_3(A)a|\le d(A)\|a\|^2$.
For PSD $A$, $d(A)\le\lambda_{\max}(A)$.
\end{proof}

\section{All PSD equality cases in order six}
A signed perfect-matching matrix is a real symmetric zero-diagonal matrix
with exactly one nonzero entry, equal to $1$ or $-1$, in every row.
It has three disjoint signed edges and satisfies $M^2=I$.

\begin{theorem}[Equality classification]\label{thm:equality}
Let $A\ne0$ be PSD of order six and let $\ell=\lambda_{\max}(A)$.
Equality in $\|\mathcal T_3(A)\|_{\op}\le\ell$ holds exactly when $A/\ell$
belongs to one of the following three families:
\begin{equation}\label{eq:families}
\boxed{ss^T/6,\qquad I-ss^T/6,\qquad (I+M)/2,}
\end{equation}
where $s\in\{\pm1\}^6$ and $M$ is any signed perfect-matching matrix.
For either $q_2(A)\le\ell^2q_1(A)/4$ or
$|\haf(A)|\le\ell q_1(A)/6$, these are all nondegenerate equality cases;
all diagonal PSD matrices additionally give equality with $q_1=q_2=0$.
\end{theorem}
\begin{proof}
Normalize $\ell=1$ and choose a unit real vector $v$ with
$v^T\mathcal T_3(A)v=\pm1$. Treat $v$ as a zero-diagonal symmetric edge array.
Let $C$ be the $6\times20$ matrix with columns $c_R(v)$ over all triples,
let $J$ exchange complementary triples, and put $P_\pm=(I\pm J)/2$.
Set
\[
H=\tfrac12CJC^T=H_+-H_-,\qquad H_\pm=\tfrac12CP_\pm C^T.
\]
The diagonal of $H$ vanishes and its off-diagonal entries are the complementary
four-variable hafnians of $v$. Since $\|C\|_F^2=4\|v\|^2=4$,
\[
\tr H_+=\tr H_-=1,\qquad v^T\mathcal T_3(A)v=\tr(AH).
\]
Equality requires $A$ to be the identity on the range of $CP_+$ and zero on
the range of $CP_-$, or conversely. These ranges are orthogonal, so
$P_+C^TCP_-=0$ and $C^TC$ commutes with $J$.

Diagonal entries of this commutation relation give the triple-balance condition
\begin{equation}\label{eq:balance}
\sum_{\{i,j\}\subset R}v_{ij}^2=
\sum_{\{i,j\}\subset R^c}v_{ij}^2\qquad (|R|=3).
\end{equation}
Its entry for the triples $\{i,j,k\}$ and $\{i,j,l\}$ gives
\[
v_{ik}v_{il}+v_{jk}v_{jl}=v_{pk}v_{pl}+v_{qk}v_{ql},
\]
where $p,q$ are the two remaining indices. Varying the two-pair split outside
$\{k,l\}$ shows that
\begin{equation}\label{eq:products}
v_{ak}v_{al}\text{ is independent of }a\notin\{k,l\}.
\end{equation}

Consider the support graph of $v$. If one vertex has neighbors $k,l$, then
\eqref{eq:products} forces both $k,l$ adjacent to all four other vertices.
Any pair among these four shares neighbor $k$, so the same condition forces
all edges among the four. A further application, using a neighbor in that
four-set, forces the edge $k,l$. Thus the graph is complete. Otherwise every
vertex has degree at most one. For any nonzero edge $ab$, applying
\eqref{eq:balance} to $\{a,b,c\}$ for each $c\notin\{a,b\}$ forces two
further disjoint edges with the same squared weight. This is a perfect
matching with constant edge modulus.

In the complete case fix $a,b$ and set $r_k=v_{ak}/v_{bk}$ for the other
four indices. Equation~\eqref{eq:products} gives $r_kr_l=1$ for every
distinct $k,l$, so all these ratios are the same sign. All incident edge
moduli, and hence all edge moduli, coincide. The row-sign ratios and symmetry
imply $v_{ij}=\alpha s_i s_j$ for $i\ne j$, with a common nonzero real
$\alpha$ and signs $s_i$.

For this dense array,
\[
H=3\alpha^2\Bigl(\prod_i s_i\Bigr)D_s(J_6-I)D_s,
\]
so $H$ is invertible and has positive and negative subspace dimensions one
and five, in either order. For a constant-modulus matching array,
\[
H=c^2\Bigl(\prod_{e\text{ in matching}}\operatorname{sign}v_e\Bigr)M,
\]
which is invertible with signature $(3,3)$. In both cases the orthogonal
ranges of $H_+$ and $H_-$ fill $\R^6$. Therefore $A$ is the positive or
negative spectral projection of $H$, giving precisely \eqref{eq:families}.

Direct substitution verifies equality in all three families, in the operator
bound and both scalar bounds. When $q_1>0$, equality in the gradient bound
forces equality in the operator bound by \eqref{eq:gradient}; equality in
the hafnian-energy bound does so through its Rayleigh inequality. Finally,
$q_1=0$ is exactly the diagonal case.
\end{proof}

\section{Scope, comparison, and a scaling correction}\label{sec:scope}
Theorem~\ref{thm:main} proves the full PSD-contraction gradient statement,
including non-idempotent inputs. It bypasses the previous nested-spectral
coefficientwise positivity approach. In particular it does not prove that
every mixed coefficient in that stronger auxiliary problem is nonnegative.
The separate problem
\[
\sup_{B\succeq0,\ \max_{x\in\{\pm1\}^6}x^TBx\le1}
|\haf(2\operatorname{offdiag}B)|=\frac1{216}
\]
remains unresolved in this work. A spectral constraint is not a Boolean-cube
constraint, and the present energy bounds alone do not close that extremum.

There is an important correction to the project's earlier operator warning.
For the conference involution $K=S/\sqrt5$, where $S^2=5I$, one has
$\|\mathcal T_3(K)\|=3/\sqrt5>1$. This refutes a unit bound for such
involutions, not the PSD bound: with $P=(I+K)/2$,
\[
\|\mathcal T_3(P)\|=\tfrac12\|\mathcal T_3(K)\|
=\frac{3}{2\sqrt5}<1.
\]
The stronger PSD operator statement previously labeled false is instead
proved by \eqref{eq:minus}--\eqref{eq:plus}.

Roos's Theorem~2.3 specializes in order six to
$q_2(A)\le3q_1(A)^2/5$. Our spectral statement uses different information:
for the three-pair projection $q_1=3/4$, it gives the exact $q_2\le3/16$,
whereas that displayed specialization gives $27/80$. This is a comparison
with a specific published bound, not a claim to supersede all hafnian
inequalities. Kummer's cited Kneser construction proves a derivative
isomorphism, not the weighted inequality here. Full-text review of the
Lindell--Staples paper was not available for this research pass; a complete
specialist priority check remains a pre-submission requirement.

A companion proof draft in the same artifact studies zero-diagonal complex
symmetric matrices. Its four-term defect identity gives all-even-dimensional
four-subhafnian bounds and quantitative phase-matching recovery, including
optimal-order stability in order six. Those complex stability statements
are complementary to the real PSD equality classification above; they
should not be described as a quantitative classification of all three PSD
families.

\section{Exact verification and reproducibility}
The analytic proofs above establish the universal statements. The software
provides independent checks with no optimization dependencies beyond NumPy.
The module \texttt{matching\_dilation} implements direct and dilation-based
operators, both PSD slack expressions, and rational six-variable certificates.
Its explicitly enumerated numerical interface is capped at order ten to
prevent accidental exponential allocations; the theorem has no dimension cap.

For orders $4,6,8,10$, tests verify the dilation on every symmetric matrix
basis element using integers. This checks all coefficients of the linear
map in those dimensions rather than only sampled matrices. Rational-instance
certificates verify $A\succeq0$, $\ell I-A\succeq0$, both operator slacks,
and the scalar corollaries by exact arithmetic, including singular pivots.
Floating-point inputs are rejected by the exact interface. Other tests cover
all projection ranks, non-idempotent matrices, sharp families, and the
conference scaling correction. The companion suite expands its quartic
identities as exact sparse polynomials and tests constructive stability.

Run the full repository suite with \texttt{python -m pytest}. For the research
overlay alone, set \texttt{PYTHONPATH=src} first. The reproducibility report
records the actual local and remote-CI environments and counts; those checks
are evidence of implementation consistency, not substitutes for the proofs.
The baseline is commit \texttt{b0143ef5052ff7caa30ce11a51e4586865b53bc3},
and the implementation is proposed in pull request 19 of
\texttt{Null-Square/no-free-collapse}.

\begin{thebibliography}{9}
\bibitem{kummer} M. Kummer, \emph{Spectral linear matrix inequalities},
Advances in Mathematics \textbf{384} (2021), 107749.
DOI: \href{https://doi.org/10.1016/j.aim.2021.107749}{10.1016/j.aim.2021.107749}.
Preprint: arXiv:2008.13452v2, especially Theorem 2.2 and Lemma 3.5.
\bibitem{roos} B. Roos, \emph{New inequalities for permanents and hafnians and
some generalizations}, Linear and Multilinear Algebra \textbf{73} (2025),
1634--1667. DOI: \href{https://doi.org/10.1080/03081087.2024.2436061}{10.1080/03081087.2024.2436061}.
Preprint: arXiv:1906.06176v2, especially Theorem 2.3.
\bibitem{lindell} T. Lindell and G. S. Staples, \emph{Norm inequalities in zeon
algebras}, Advances in Applied Clifford Algebras \textbf{29} (2019), article 13.
DOI: \href{https://doi.org/10.1007/s00006-018-0934-z}{10.1007/s00006-018-0934-z}.
\end{thebibliography}
\end{document}
